\section{Results}\label{results}

\subsection{Overall Performance}\label{overall_performance}

\autoref{tab:performance} summarises the final test performance of the pretrained baseline, SFT and
OPSD conditions across all three VLMs. BertScore-F1 is reported as the primary metric with
normalised EM as the secondary metric. Best result per model and metric is shown in bold:

\begin{table}[h!]
\centering
\caption{Final test performance across pretrained, SFT, and OPSD conditions.}
\label{tab:performance}
\resizebox{\columnwidth}{!}{
\begin{tabular}{lcccccc}
\hline
\textbf{Model} &
\textbf{Base F1} &
\textbf{SFT F1} &
\textbf{OPSD F1} &
\textbf{Base EM} &
\textbf{SFT EM} &
\textbf{OPSD EM} \\
\hline

SmolVLM &
0.9256 &
\textbf{0.9712} &
0.9259 &
0.5651 &
\textbf{0.6914} &
0.5690 \\

BLIP &
0.9569 &
\textbf{0.9727} &
0.9570 &
0.5833 &
\textbf{0.7057} &
0.5833 \\

MiniCPM &
0.9793 &
\textbf{0.9806} &
0.9805 &
0.7904 &
0.7904 &
\textbf{0.7956} \\

\hline
\end{tabular}
}
\end{table}

For SmolVLM, SFT produced a substantial improvement over the pretrained baseline, increasing
BERTScore-F1 from 0.9256 to 0.9712 and exact match from 0.5651 to 0.6914. OPSD, in contrast,
achieved 0.9259 BERTScore-F1 and 0.5690 exact match, remaining close to the pretrained baseline.
A similar pattern was observed for BLIP. SFT improved BERTScore-F1 from 0.9569 to 0.9727 
and exact match from 0.5833 to 0.7057. OPSD achieved a BERTScore-F1 of 0.9570 and an 
exact match of 0.5833, again closely matching the pretrained model. The difference was 
considerably smaller for MiniCPM. SFT achieved a BERTScore-F1 of 0.9806, while 
OPSD achieved 0.9805. Exact-match performance was slightly higher under OPSD, 
increasing from 0.7904 under SFT to 0.7956.

The effectiveness of the privileged teacher was also assessed before fine-tuning by comparing
the crop-conditioned and full-image predictions of each pretrained model. The mean BERTScore-F1
advantage of the crop-conditioned teacher was $-0.00067$ for SmolVLM, $+0.00188$ for BLIP
and $-0.00027$ for MiniCPM. These results indicate that the privileged crops did not
provide a consistent average performance advantage across the three architectures.

\subsection{Paired Statistical Comparison}

To directly compare SFT and OPSD, a paired analysis was performed over the same held-out test
set. \autoref{tab:paired_differences} reports the mean per-example BERTScore-F1 difference together with
the corresponding 95\% bootstrap confidence interval:

\begin{table}[h!]
\centering
\caption{Paired OPSD--SFT BERTScore-F1 differences on the test set.}
\label{tab:paired_differences}
\begin{tabular}{lcc}
\hline
\textbf{Model} &
$\boldsymbol{\Delta_{\mathrm{OPSD-SFT}}}$ &
\textbf{95\% CI} \\
\hline

SmolVLM &
$-0.0453$ &
$[-0.0483, -0.0423]$ \\

BLIP &
$-0.0157$ &
$[-0.0192, -0.0123]$ \\

MiniCPM &
$-0.0001$ &
$[-0.0026, 0.0024]$ \\

\hline
\end{tabular}
\end{table}

For SmolVLM and BLIP, the confidence intervals are entirely below zero, indicating a 
consistent advantage for SFT over OPSD on the held-out test set. SmolVLM showed the 
largest difference, with OPSD trailing SFT by 0.0453 BERTScore-F1, followed by BLIP 
with a difference of 0.0157. For MiniCPM, the difference between OPSD and SFT was 
negligible (-0.0001), and the confidence interval crossed zero. Although OPSD achieved a small
exact-match improvement of approximately 0.52 percentage points, the BERTScore-F1 results do not
establish a statistically clear advantage for either method.

\subsection{Per-example Win/Loss Analysis}

To complement the aggregate BertScore-F1 comparison, the paired test predictions were
also analysed at the individual-question level. For each test example, OPSD was classified
as better than SFT when its per-example BertScore-F1 was higher, SFT was classified as better
when its score was higher and equal scores were treated as ties. \autoref{tab:per_example_comparison}
reports the resulting proportions across the 768 paired test examples:

\begin{table}[h!]
\centering
\caption{Per-example BERTScore-F1 comparison between OPSD and SFT.}
\label{tab:per_example_comparison}
\resizebox{\columnwidth}{!}{
\begin{tabular}{lccc}
\hline
\textbf{Model} &
\textbf{OPSD Better} &
\textbf{SFT Better} &
\textbf{Tie} \\
\hline

SmolVLM &
18.23\% &
81.77\% &
0.00\% \\

BLIP &
5.73\% &
21.22\% &
73.05\% \\

MiniCPM &
8.33\% &
8.98\% &
82.68\% \\

\hline
\end{tabular}
}
\end{table}

The per-example analysis reinforces the aggregate findings for SmolVLM. SFT achieved a 
higher BERTScore-F1 on 81.77\% of test questions, compared with 18.23\% for OPSD, with 
no ties observed. This is consistent with the comparatively large negative paired 
difference reported in Table~\ref{tab:paired_differences}. For BLIP, the majority of examples produced 
identical BERTScore-F1 values under the two post-training strategies, with 73.05\% of 
examples resulting in ties. Among the remaining examples, SFT was better on 21.22\% of 
the test set, while OPSD was better on 5.73\%. This indicates that the aggregate SFT 
advantage was driven by a relatively smaller subset of examples rather than by 
differences across the entire test set. MiniCPM exhibited the highest degree of 
agreement between SFT and OPSD, with 82.68\% of test examples resulting in ties. 
OPSD achieved the higher BERTScore-F1 on 8.33\% of examples, while SFT was better on 
8.98\%. The near-symmetric win rates are consistent with the negligible aggregate 
difference between the two methods and the bootstrap confidence interval crossing zero.

Overall, the results for RQ1 show that OPSD did not outperform conventional SFT in 
terms of BERTScore-F1. SFT was substantially stronger for SmolVLM and BLIP, while 
the two methods achieved near-identical performance for MiniCPM. With respect to RQ2, 
the relative OPSD--SFT gap varied considerably across architectures. 
The difference decreased from (-0.0453) for SmolVLM to (-0.0157) for BLIP and to 
approximately zero for MiniCPM. This findings indicate that the relative effectiveness
of OPSD varied across the evaluated VLMs, suggesting that model characteristics may influence the
benefits of privileged-context self-distillation.

